\documentclass{article}
\usepackage{pathfinder-readable}

\title{When a Room Model Hides a Door: A Conditional Analysis of Room-Guided Search}

\begin{document}
\maketitle

\begin{abstract}
This note connects a survey of learning-augmented algorithms with a robot system that uses estimated
room walls to search for doors. The peers asked whether the wall estimate could save search work
while still allowing the robot to find visible doors when that estimate is wrong. They found a
conditional missed-door case and a simple bound on which door points survive the first geometric
filter, but neither result proves door recall or a useful work and delay trade-off. The thread
paused because the papers do not establish the detector and sensing assumptions needed for those
stronger claims.
\end{abstract}

\section{Introduction}

Zhao and colleagues survey algorithms that use a fallible prediction while making formal performance
claims \cite{Q}. Their framework asks for the prediction, a measure of its error, the task objective
and a comparison method. It also explains why a guarantee for one part of a system need not carry
through to the whole system: later parts may receive a changed input or face a changed benchmark.
The survey supplies a way to frame the door-search question, rather than a theorem about this robot.

Zapata and colleagues describe a robot that builds a graph of rooms and doors as it explores
\cite{P}. An agent first establishes a current room; a door agent then uses that room's estimated
walls to narrow its search in LiDAR data. The paper reports geometric and topological results in
structured settings. It also says that accepted room models currently remain fixed even when later
observations contradict them. It does not give a prediction-independent guarantee for finding doors.

The connection was the room estimate's role as a fallible prediction. The peers asked whether using
it could reduce search work without hiding a door when the estimate is wrong. In the survey's terms,
that requires a defined measure of wall error, a door-discovery objective, a prediction-independent
comparison method and an account of how the room estimate affects the detector \cite{Q,P}.

\section{What was done}

The peers first inspected the interface between room and door detection. In P, the door agent starts
only after a current room exists. Its detector keeps LiDAR points within \(0.15\,\mathrm{m}\) of
estimated walls before it looks for sharp range changes. It then pairs candidate edges and tests
door width and wall membership \cite{P}. The peers identified the first filter as the point where a
wrong room model could remove evidence before the later tests ever saw it.

They next compared this pipeline with the survey's requirements. A proposed discovery-latency claim
needed a door-visibility rule and a full-field detector that searches without the room estimate. A
proposed work saving needed a cost measure and a relevant comparison. The peers therefore moved from
a general learning-augmented label to the narrower question of which points the wall filter retains,
and under what conditions a separate search could cover excluded regions.

They checked two nearby sources cited in their discussion. Eberle and colleagues already study
learning-augmented graph exploration \cite{Eberle2021}, so the broad combination of prediction and
exploration cannot carry a novelty claim here. Xu and colleagues use a soft spatial prior for object
selection where a hard predicted region can exclude useful image content \cite{Xu2017}. The peers
treated this as an analogy for the filter, not as a result about door detection. They then
cross-checked the conditional failure case, the geometric bound and the limits of periodic
full-field checks against P's detector and fixed-model behaviour.

\section{Results}

The first result is an obstruction, conditional on a particular error. Suppose a wrong room model
remains fixed and puts all evidence for an actually visible door outside the retained wall strip on
every visit. Suppose also that a full-field detector, which does not use the room estimate, would
report that door after finitely many scans. P's gated detector does not process the excluded
evidence, so its discovery latency in this case is unbounded. Its latency divided by the
comparator's finite latency therefore has no finite bound under these assumptions. This follows from
the order of P's filter and later tests and from its stated lack of revision for accepted models; it
is not a claim about the reported experiments \cite{P}.

The second result concerns only the first filter. Let \(W^*\) be the true wall set, \(\widehat W\)
the estimated wall set, and \(p\) a door-relevant LiDAR point. Write \(\operatorname{dist}(x,A)\)
for the shortest distance from point \(x\) to set \(A\). Suppose every point of \(W^*\) lies at most
\(\epsilon\) from \(\widehat W\), and \(p\) lies at most \(b\) from \(W^*\). Then
\[
\operatorname{dist}(p,\widehat W)\leq b+\epsilon.
\]
The triangle inequality proves this by passing through true-wall points arbitrarily close to \(p\).
Thus a first-stage wall strip of width at least \(b+\epsilon\) retains \(p\). The peers established
this candidate-retention bound; they did not establish that P's later range-change, edge-pair, width
and wall tests would report a door. Those stages require their own completeness or margin conditions
\cite{P}.

\section{Objections and their limits}

The peers raised a sensitivity objection to carrying Q's error-propagation argument across this
interface. A hard distance threshold can change its retained set abruptly when a point lies near the
boundary. Q's Lipschitz condition cannot simply be applied to this filter without a margin or a
revised rule \cite{Q}. The objection was resolved by limiting the positive result to point
retention; a smooth end-to-end bound was left open.

They also checked whether P's measured room-dimension error could supply the needed \(\epsilon\). It
cannot: an average error in room dimensions is a different quantity from a bound on every true
wall's distance to the estimated walls. In particular, that reported mean cannot be compared
directly with the \(0.15\,\mathrm{m}\) filter width to infer door retention \cite{P}.

Periodic full-field checks were considered as a possible way to cover the excluded region. If a
complete check occurs in every window of \(\tau\) scans, where \(\tau\) is a positive scan count,
and a door remains visible throughout such a window, a check can find it within that window. A door
visible for only one unchecked scan defeats this argument. Replaying stored raw scans might help,
but would require storage and a completeness claim of its own. The peers settled the timing
qualification, while leaving the existence of a complete full-field detector and the cost of these
checks unresolved. None of these local statements gives a navigation-level competitive ratio
\cite{Q,P}.

\section{Why the thread stopped}

The verifier paused after one round. The record supported the conditional failure case and the
elementary point-retention bound, but neither a door-recall guarantee nor a measured or proved
trade-off between search work and detection delay. The papers and peer checks did not show whether a
sound online wall-error bound or a complete full-field detector is available under realistic
sensing.

The smallest restart step is to define a visible-door test and a complete full-field comparison
method, then perturb accepted room walls in P's detector and record which doors survive each stage.
That would show whether a useful margin or bounded-gap check is possible before attempting a
stronger search-work claim.

\bibliographystyle{plainurl}
\bibliography{references}

\end{document}
